\documentclass[10pt,a4paper]{amsart}
\usepackage{amsmath,amssymb,amsthm,graphicx}
% Retained preamble for the Stacks Project chapter snapshot in this directory.
% The chapter's own preamble.tex is a copy of the shared Stacks Project preamble; the block
% below is sliced MECHANICALLY, line for line, from it. Nothing was added, removed or
% redefined; the wrapper class, the xr-hyper/hyperref wiring and the externaldocument list
% are omitted because the assembler supplies the document class.
%
% lines taken from preamble.tex: 3, 5, 9, 11, 12, 156, 157, 158, ... (51 lines total)
%
% package coverage restored from the chapter's own preamble (2026-10-13): xy, xy-2cell

\usepackage[all]{xy}
\xyoption{2cell}
\UseAllTwocells

\newenvironment{reference}{\comment}{\endcomment}
\newenvironment{slogan}{\comment}{\endcomment}
\newenvironment{history}{\comment}{\endcomment}
\theoremstyle{plain}
\newtheorem{theorem}[subsection]{Theorem}
\newtheorem{proposition}[subsection]{Proposition}
\newtheorem{lemma}[subsection]{Lemma}
\theoremstyle{definition}
\newtheorem{definition}[subsection]{Definition}
\newtheorem{example}[subsection]{Example}
\newtheorem{exercise}[subsection]{Exercise}
\newtheorem{situation}[subsection]{Situation}
\theoremstyle{remark}
\newtheorem{remark}[subsection]{Remark}
\newtheorem{remarks}[subsection]{Remarks}
\numberwithin{equation}{subsection}
\def\lim{\mathop{\mathrm{lim}}\nolimits}
\def\colim{\mathop{\mathrm{colim}}\nolimits}
\def\Spec{\mathop{\mathrm{Spec}}}
\def\Hom{\mathop{\mathrm{Hom}}\nolimits}
\def\Ext{\mathop{\mathrm{Ext}}\nolimits}
\def\SheafHom{\mathop{\mathcal{H}\!\mathit{om}}\nolimits}
\def\SheafExt{\mathop{\mathcal{E}\!\mathit{xt}}\nolimits}
\def\Sch{\mathit{Sch}}
\def\Mor{\mathop{\mathrm{Mor}}\nolimits}
\def\Ob{\mathop{\mathrm{Ob}}\nolimits}
\def\Sh{\mathop{\mathit{Sh}}\nolimits}
\def\NL{\mathop{N\!L}\nolimits}
\def\CH{\mathop{\mathrm{CH}}\nolimits}
\def\proetale{{pro\text{-}\acute{e}tale}}
\def\etale{{\acute{e}tale}}
\def\QCoh{\mathit{QCoh}}
\def\Ker{\mathop{\mathrm{Ker}}}
\def\Im{\mathop{\mathrm{Im}}}
\def\Coker{\mathop{\mathrm{Coker}}}
\def\Coim{\mathop{\mathrm{Coim}}}
\DeclareMathSymbol{\boxtimes}{\mathbin}{AMSa}{"02}
\def\QCohstack{\mathcal{QC}\!\mathit{oh}}
\def\Cohstack{\mathcal{C}\!\mathit{oh}}
\def\Spacesstack{\mathcal{S}\!\mathit{paces}}
\def\Quotfunctor{\mathrm{Quot}}
\def\Hilbfunctor{\mathrm{Hilb}}
\def\Curvesstack{\mathcal{C}\!\mathit{urves}}
\def\Polarizedstack{\mathcal{P}\!\mathit{olarized}}
\def\Complexesstack{\mathcal{C}\!\mathit{omplexes}}
\def\Pic{\mathop{\mathrm{Pic}}\nolimits}
\def\Picardstack{\mathcal{P}\!\mathit{ic}}
\def\Picardfunctor{\mathrm{Pic}}
\def\Deformationcategory{\mathcal{D}\!\mathit{ef}}
\usepackage{cleveref}
\crefname{originallemma}{original lemma}{original lemmas}
\Crefname{originallemma}{Original lemma}{Original lemmas}
\makeatletter
\newlabel{equation-to-solve-ringed-spaces}{{08U7}{source}{}{sourcepointer-1}{}}
\newlabel{equation-to-solve-ringed-spaces@cref}{{[originallemma][0][]08U7}{[1][0][]source}{}{}{}}
\makeatother
\crefname{originallemma}{original lemma}{original lemmas}
\Crefname{originallemma}{Original lemma}{Original lemmas}
\makeatletter
\newlabel{lemma-NL-represent-ext-class}{{08UB}{source}{}{sourcepointer-2}{}}
\newlabel{lemma-NL-represent-ext-class@cref}{{[originallemma][0][]08UB}{[1][0][]source}{}{}{}}
\makeatother
\crefname{originallemma}{original lemma}{original lemmas}
\Crefname{originallemma}{Original lemma}{Original lemmas}
\makeatletter
\newlabel{lemma-huge-diagram-ringed-spaces}{{08U8}{source}{}{sourcepointer-3}{}}
\newlabel{lemma-huge-diagram-ringed-spaces@cref}{{[originallemma][0][]08U8}{[1][0][]source}{}{}{}}
\makeatother
\title{Stacks, Deformation Theory: lemma-choices-ringed-spaces}
\author{The Stacks Project Authors}
\date{stacks-project:defos.tex:local-snapshot-20261012; GNU FDL 1.2 or later}
\begin{document}
\maketitle
\noindent\textit{Editorial note.} Statement span is lines 2216-2221 of sources/stacks-defos/pinned-original-source.tex; proof span is lines 2223-2278, the complete author proof (54 body lines).\par\medskip
\noindent\textit{Editorial note.} Navigation stubs: 3 cited-but-not-included labels are supplied as non-load-bearing pointers with their pinned tags where the snapshot tag index has them.\par\medskip
\noindent\textit{Editorial note.} extra\_packages is empty on purpose: the retained preamble already loads xy with options and a second declaration clashes in the assembler wrapper.\par\medskip
\noindent\textit{Editorial note.} Text is retained verbatim, including the author's own typos and placeholders.\par\medskip
\noindent\textit{Editorial note.} Queue index 12 of queue\_defos.json (pb=2411, stmt 2216-2221, proof 2223-2278).\par\medskip
\noindent\textit{Allowed standard prerequisites (editorial).}\begin{itemize}
\item Explicit prerequisite boundary. The cited inputs are prerequisites; the argument in the retained span is not replaced by any cited result.
\end{itemize}\medskip
\ifdefined\hypertarget\hypertarget{sourcepointer-1}{}\fi
\noindent\textit{Source reference.} Original lemma 08U7: sources/stacks-defos/pinned-original-source.tex, source lines 1935--1943. Cited by the retained statement or proof but not part of the collected spans; supplied only as a navigation pointer. The argument does not depend on its content.\par\medskip
\ifdefined\hypertarget\hypertarget{sourcepointer-2}{}\fi
\noindent\textit{Source reference.} Original lemma 08UB: sources/stacks-defos/pinned-original-source.tex, source lines 2158--2167. Cited by the retained statement or proof but not part of the collected spans; supplied only as a navigation pointer. The argument does not depend on its content.\par\medskip
\ifdefined\hypertarget\hypertarget{sourcepointer-3}{}\fi
\noindent\textit{Source reference.} Original lemma 08U8: sources/stacks-defos/pinned-original-source.tex, source lines 1956--1972. Cited by the retained statement or proof but not part of the collected spans; supplied only as a navigation pointer. The argument does not depend on its content.\par\medskip
\begin{lemma}
\label{lemma-choices-ringed-spaces}
If there exists a solution to (\ref{equation-to-solve-ringed-spaces}),
then the set of isomorphism classes of solutions is principal homogeneous
under $\Ext^1_{\mathcal{O}_X}(\NL_{X/S}, \mathcal{G})$.
\end{lemma}

\begin{proof}
We observe right away that given two solutions $X'_1$ and $X'_2$
to (\ref{equation-to-solve-ringed-spaces}) we obtain by
Lemma \ref{lemma-huge-diagram-ringed-spaces} an obstruction element
$o(X'_1, X'_2) \in \Ext^1_{\mathcal{O}_X}(\NL_{X/S}, \mathcal{G})$
to the existence of a map $X'_1 \to X'_2$. Clearly, this element
is the obstruction to the existence of an isomorphism, hence separates
the isomorphism classes. To finish the proof it therefore suffices to
show that given a solution $X'$ and an element
$\xi \in \Ext^1_{\mathcal{O}_X}(\NL_{X/S}, \mathcal{G})$
we can find a second solution $X'_\xi$ such that
$o(X', X'_\xi) = \xi$.

\medskip\noindent
Pick $\alpha : \mathcal{E} \to \mathcal{O}_X$ as in
Lemma \ref{lemma-NL-represent-ext-class}
for the class $\xi$. Consider the surjection
$f^{-1}\mathcal{O}_S[\mathcal{E}] \to \mathcal{O}_X$
with kernel $\mathcal{I}$ and corresponding naive cotangent complex
$\NL(\alpha) = (\mathcal{I}/\mathcal{I}^2 \to
\Omega_{f^{-1}\mathcal{O}_S[\mathcal{E}]/f^{-1}\mathcal{O}_S}
\otimes_{f^{-1}\mathcal{O}_S[\mathcal{E}]} \mathcal{O}_X)$.
By the lemma $\xi$ is the class of a morphism
$\delta : \mathcal{I}/\mathcal{I}^2 \to \mathcal{G}$.
After replacing $\mathcal{E}$ by
$\mathcal{E} \times_{\mathcal{O}_X} \mathcal{O}_{X'}$ we may also assume
that $\alpha$ factors through a map
$\alpha' : \mathcal{E} \to \mathcal{O}_{X'}$.

\medskip\noindent
These choices determine an $f^{-1}\mathcal{O}_{S'}$-algebra map
$\varphi : f^{-1}\mathcal{O}_{S'}[\mathcal{E}] \to \mathcal{O}_{X'}$. Let
$\mathcal{I}' = \Ker(f^{-1}\mathcal{O}_{S'}[\mathcal{E}] \to \mathcal{O}_X)$.
Observe that $\varphi$ induces a map
$\varphi|_{\mathcal{I}'} : \mathcal{I}' \to \mathcal{G}$
and that $\mathcal{O}_{X'}$ is the pushout, as in the following
diagram
$$
\xymatrix{
0 \ar[r] & \mathcal{G} \ar[r] & \mathcal{O}_{X'} \ar[r] &
\mathcal{O}_X \ar[r] & 0 \\
0 \ar[r] & \mathcal{I}' \ar[u]^{\varphi|_{\mathcal{I}'}} \ar[r] &
f^{-1}\mathcal{O}_{S'}[\mathcal{E}] \ar[u] \ar[r] &
\mathcal{O}_X \ar[u]_{=} \ar[r] & 0
}
$$
Let $\psi : \mathcal{I}' \to \mathcal{G}$ be the sum of the map
$\varphi|_{\mathcal{I}'}$ and the composition
$$
\mathcal{I}' \to \mathcal{I}'/(\mathcal{I}')^2 \to
\mathcal{I}/\mathcal{I}^2 \xrightarrow{\delta} \mathcal{G}.
$$
Then the pushout along $\psi$ is an other ring extension
$\mathcal{O}_{X'_\xi}$ fitting into a diagram as above.
A calculation (omitted) shows that $o(X', X'_\xi) = \xi$ as desired.
\end{proof}

\end{document}
